%!TEX root = ./poly.tex
%
\section{\tops{$\RO(\T)$}{RO(T)}-graded \tops{$\TP$}{TP} and \tops{$\TC^-$}{TC-} of quasiregular semiperfectoid rings}
\label{sec:qrsp}


\newcommand{\vprism}{\prismcpl^{(1)}}

In this section we study the $\RO(\T)$ (really $\RO(C_{p^\infty})$)-graded $\TCmin$ and $\TP$ of quasiregular semiperfectoid rings $S$. We compute the relevant groups in \sec\ref{sub:qrsp-fractional}; the result is essentially the multi-Rees algebra of the Nygaard and $\phi^\bullet(I)$-adic filtrations on $\prismcpl_S[1/\phi^\bullet(I)]$, so that we encounter \emph{prismatic fractional ideals} like $I^2 \tnsr \phi(I)^{-3}$. We also give integral generalizations of some of our results to even $\E_\infty$-rings with $\T$-action. These computations ignore the distinction between $\lambda^n$ and $\lambda^{p^{v_p(n)}}$; in \sec\ref{sub:integral} we explain how to put this distinction back in, which does not affect module structures but \emph{does} twist restriction ($F$) maps, and becomes important when studying norms. Finally, in \sec\ref{sub:iterated-nygaard} we extend our computations to the spectra $\TC^{-n} \defeq (\TR^n)^{h\T/C_{p^n}}$.

\subsection{Prismatic fractional ideals}
\label{sub:qrsp-fractional}
In this section we explain how to see Theorem \ref{thm:qrsp}. We fix a $\delta$-ring $A$ with spherical lift $\S_A$, a relatively quasiregular semiperfectoid $A$-algebra $S$, and abbreviate $\TCmin\defeq\TCmin(S\rel\S_A)$, $\TP\defeq\TP(S\rel\S_A)$, $\vprism \defeq\vprism_{S/A}$, etc.

Our main tools will be the cell structure
\begin{equation}
  \label{eq:cell}
  \T/C_{n+} \to S^0 \to S^{\lambda^n}
\end{equation}
and the motivic filtrations
\begin{align*}
  \TCmin_{2i} &= \Nyg^{\ge i}\vprism\{i\} & \TP_{2i} &= \vprism\{i\}\\
  \THH^{hC_{p^n}}_{2i} &= \frac{\Nyg^{\ge i}\vprism\{i\}}{I_n\tnsr\Nyg^{\ge i+1}\vprism\{i\}} & \THH^{tC_{p^n}}_{2i} &= \frac{\vprism\{i\}}{I_n\vprism\{i\}}
\end{align*}
from \cite[Theorem 1.12(4)]{BMS2} and \cite[Corollary 3.3]{RiggenbachTruncated}. We will just explain the case of $\TCmin$, leaving the (easier) case of $\TP$ to the reader.

\begin{remark}
  \label{rmk:no-extns}
  Since $\THH$ of perfectoid rings is even, there are invertible classes $u_{\lambda_k}\in\TCmin_{2-\lambda_k}$ giving non-canonical isomorphisms
  \[ \TCmin_\alpha \cong \TCmin_{2d_0(\alpha)} \]
  for every $\alpha\in\RU(\T)$ \cite[Proposition 4.2]{RiggenbachDoublePoints}. Of course the whole point of this section is to avoid these classes and instead give canonical identifications. However, this observation is useful for solving extensions and identifying images.
\end{remark}

\begin{example}
  \label{ex:qrsp-pos}
  The cell structure \eqref{eq:cell} gives a short exact sequence
  \[\xymatrix{
      0 \ar[r] & \TCmin_{2i+\lambda_k} \ar[r]^-{a_{\lambda_k}} \ar@{=}[d] & \TCmin_{2i} \ar[r] \ar@{=}[d] & \THH^{hC_{p^k}}_{2i} \ar[r] \ar@{=}[d] & 0\\
      0 \ar[r] & ? \ar[r] & \Nyg^{\ge i}\vprism\{i\} \ar[r] & \dfrac{\Nyg^{\ge i}\vprism\{i\}}{I_k\tnsr\Nyg^{\ge i+1}\vprism\{i\}} \ar[r] & 0
  }\]
  which shows that $\TCmin_{2i+\lambda_k}=I_k\tnsr\Nyg^{\ge i+1}\vprism\{i\}$, as claimed by Theorem \ref{thm:qrsp}.
\end{example}

\begin{example}
  The cell structure \eqref{eq:cell} gives a short exact sequence
  \[\xymatrix{
      0 \ar[r] & \TCmin_{2i} \ar[r]^-{a_{\lambda_k}} \ar@{=}[d] & \TCmin_{2i-\lambda_k} \ar[r] \ar@{=}[d] & \THH^{hC_{p^k}}_{2(i-1)} \ar[r] \ar@{=}[d] & 0\\
      0 \ar[r] & \Nyg^{\ge i}\vprism\{i\} \ar[r] & ? \ar[r] & \dfrac{\Nyg^{\ge i-1}\vprism\{i-1\}}{I_k\tnsr\Nyg^{\ge i}\vprism\{i-1\}} \ar[r] & 0
  }\]
  Rewriting the final term as
  \[
    \dfrac{\Nyg^{\ge i-1}\vprism\{i-1\}}{I_k\tnsr\Nyg^{\ge i}\vprism\{i-1\}} 
    =
    \dfrac{I_k^{-1}\tnsr\Nyg^{\ge i-1}\vprism\{i\}}{\Nyg^{\ge i}\vprism\{i\}} 
  \]
  and using Remark \ref{rmk:no-extns}, we see that $\TCmin_{2i-\lambda_k} = I_k^{-1}\tnsr\Nyg^{\ge i-1}\vprism\{i\}$, as claimed by Theorem \ref{thm:qrsp}.

  Alternatively, we can use the multiplication map
  \[\xymatrix{
      \TCmin_{2i-\lambda_k} \tnsr \TCmin_{\lambda_k-2} \ar[r] \ar@{=}[d] & \TCmin_{2(i-1)} \ar@{=}[d]\\
      \TCmin_{2i-\lambda_k} \tnsr I_k\tnsr\vprism\{-1\} \ar[r] & \Nyg^{\ge i-1}\vprism\{i-1\}
  }\]
  where the left vertical identification is Example \ref{ex:qrsp-pos}. Using Remark \ref{rmk:no-extns}, we see that the rows are isomorphisms. We thus get
  \[ \TCmin_{2i-\lambda_k} = \Nyg^{\ge i-1}\vprism\{i-1\}\tnsr\left(I_k\tnsr\vprism\{-1\}\right)^{-1} = I_k^{-1}\Nyg^{\ge i-1}\vprism\{i\}. \]
\end{example}

\begin{example}
  We now explain how to mix representations. If $k\le n$, then $\res^\T_{C_{p^k}} \lambda_n$ is trivial, but $\res^\T_{C_{p^n}} \lambda_k$ is not unless $k=n$. We will thus induct by using $\lambda_n$ before we use $\lambda_k$. Now the cell structure \eqref{eq:cell} gives a short exact sequence
  \[\xymatrix{
      0 \ar[r] & \TCmin_{2i+\lambda_k-\lambda_n} \ar[r]^-{a_{\lambda_k}} \ar@{=}[d] & \TCmin_{2i-\lambda_n} \ar[r] \ar@{=}[d] & \THH^{hC_{p^k}}_{2(i-1)} \ar[r] \ar@{=}[d] & 0\\
      0 \ar[r] & ? \ar[r] & I_n^{-1}\Nyg^{\ge i-1}\vprism\{i\} \ar[r] & \dfrac{\Nyg^{\ge i-1}\vprism\{i-1\}}{I_k\tnsr\Nyg^{\ge i}\vprism\{i-1\}} \ar[r] & 0.
  }\]
  Rewriting the final term as
  \[
    \dfrac{\Nyg^{\ge i-1}\vprism\{i-1\}}{I_k\tnsr\Nyg^{\ge i}\vprism\{i-1\}}
    =
    \dfrac{I_k^{-1}\tnsr\Nyg^{\ge i-1}\vprism\{i\}}{\Nyg^{\ge i}\vprism\{i\}}
    =
    \dfrac{I_n^{-1}\tnsr\Nyg^{\ge i-1}\vprism\{i\}}{I_k\tnsr I_n^{-1}\tnsr\Nyg^{\ge i}\vprism\{i\}},
  \]
  we see that $\TCmin_{2i+\lambda_k-\lambda_n}=I_k\tnsr I_n^{-1}\tnsr\Nyg^{\ge i}\vprism\{i\}$, as claimed by Theorem \ref{thm:qrsp}.

  Alternatively, we can use the multiplication map
  \[\xymatrix{
      \TCmin_{2i+\lambda_k-\lambda_n} \tnsr \TCmin_{\lambda_n-2} \ar[r] \ar@{=}[d] & \TCmin_{2(i-1)+\lambda_k} \ar@{=}[d]\\
      \TCmin_{2i+\lambda_k-\lambda_n} \tnsr I_n\tnsr\vprism\{-1\} \ar[r] & I_k\tnsr\Nyg^{\ge i}\vprism\{i-1\}
  }\]
  where the left vertical identification is Example \ref{ex:qrsp-pos}. Using Remark \ref{rmk:no-extns}, we see that the rows are isomorphisms. We thus get
  \[ \TCmin_{2i+\lambda_k-\lambda_n} = I_k\tnsr\Nyg^{\ge i}\vprism\{i-1\}\tnsr\left(I_n\tnsr\vprism\{-1\}\right)^{-1} = I_k\tnsr I_n^{-1}\tnsr\Nyg^{\ge i}\vprism\{i\}. \]
\end{example}

Combining the techniques in the above examples, the general case of Theorem \ref{thm:qrsp} follows immediately.

\begin{remark}
  \label{rmk:even}
  The only inputs needed for Theorem \ref{thm:qrsp} were:
  \begin{itemize}
    \item $\THH(S\rel\S_A)$ is orientable in the Borel category \cite[Proposition 4.2]{RiggenbachDoublePoints};
    \item a description of the Mackey functors $\mpi_\rog\THH(S\rel\S_A)^h$ and $\mpi_\rog\THH(S\rel\S_A)^t$ \cite[Corollary 3.3]{RiggenbachTruncated}.
  \end{itemize}
  Riggenbach in fact shows the first point for all $\T$-equivariant even $\E_\infty$-rings, while he and Antieau provide a version of the second in this generality in \cite[Corollary 3.20]{CycSyn}. We now explain this integral generalization (Theorem \ref{thm:even}).

  Let $E$ be an even $\E_\infty$-ring with $\T$-action. Set $A = \pi_0 E^{t\T}$, $A\{1\} = \pi_2 E^{t\T}$, and $A\{n\} = A\{1\}^{\tnsr n}$; by the Tate spectral sequence, we have
  \[ \pi_* E^{t\T} = \bigoplus_{i\in\Z} A\{i\}[2i] \]

  Let $\Fil^\bullet A\{i\}$ be the abutment filtration of the Tate spectral sequence. By the homotopy fixed points spectral sequence, we have
  \[ \pi_* E^{h\T} = \bigoplus_{i\in \Z} \Fil^i A\{i\}[2i] \]
  The natural map
  \[
    \Fil^i A\{i\} \tnsr_A A\{-i\} \to \Fil^i A
  \]
  induces an isomorphism
  \[ \pi_{2i} E = \gr^i A\{i\} \iso \gr^i A \]
  and therefore the short exact sequences
  \[ 0 \to E^{h\T}_{2i+(j+1)\lambda}\xra{a_\lambda} E^{h\T}_{2i+j\lambda} \to \pi_{2(i+j)} E \to 0 \]
  give identifications
  \[ E^{h\T}_{2i+j\lambda} = \Fil^{i+j}A\{i\} \]
  for all $i,j\in\Z$, as before.

  To go further, let $\G$ be the formal group law of $E$, and let $t\in A\{-1\}$ be a choice of complex orientation. Write $[n]_\G(t)\in A\{-1\}$ for the $n$-series of $\G$, $\<n\>_\G(t) = \frac{[n]_\G(t)}{t} \in A$ for the divided $n$-series; we sometimes suppress the dependence on $t$ and just write $\<n\>_\G$ (since we only care about the ideal $\<n\>_\G(t)A$). Assume furthermore that $\<n\>_\G(t)$ is a non-zerodivisor for all $n$. By \cite[Corollary 3.20]{CycSyn}, we have
  \begin{align*}
    \pi_{2i} E^{hC_n} &= \frac{\Fil^i A\{i\}}{\<n\>_\G\Fil^{i+1}A\{i\}} &
    \pi_{2i}E^{tC_n} &= \frac{A\{i\}}{\<n\>_\G}
  \end{align*}
  It follows that
  \begin{align*}
    E^{h\T}_{2i+j\lambda^n} &= \Fil^{i+j} A\{i\} \tnsr_A (\<n\>_\G A)^{\tnsr j}\\
    E^{t\T}_{2i+j\lambda^n} &= A\{i\} \tnsr_A (\<n\>_\G A)^{\tnsr j}
  \end{align*}

  To go fully general, we need a way to describe the geometric fixed points. Recall $\Phi_{d,\G}(t)$ from Definition \ref{def:fgl-cyclotomic}. In our topological situation, we have $\Phi_{d,\G}(t) \in A$ for $d>1$, and we suppress the dependence on $t$ and simply write $\Phi_{d,\G}\in A$. Note that $\Phi_{n,\G}=0$ in $\pi_0 E^{\tau C_n}$ (proper Tate construction), and heuristically $A\{1\} = \bigotimes_{n>1} \Phi_{n,\G}A$. With these notations, Theorem \ref{thm:even} follows immediately.
\end{remark}

\begin{example}
  Let $E$ be an even $\E_\infty$-ring with $\T$-action, and assume that $\<n\>_\G \in A=\pi_0 E^{t\T}$ is a non-zerodivisor for all $n$. Then for all $i\in\Z$,
  \begin{align*}
    E^{h\T}_{(n)_{\lambda^m} + 2i}
      &= \Fil^{i+n} A\{i+1\} \tnsr \<n\>_{\G^{(m)}}? &
    E^{t\T}_{(n)_{\lambda^m} + 2i}
      &= A\{1+i\} \tnsr \<n\>_{\G^{(m)}}?\\
    E^{h\T}_{\lambda^m(n)_{\lambda^m} + 2i}
      &= \Fil^{i+n} A\{i\} \tnsr \<n\>_{\G^{(m)}}! &
    E^{t\T}_{\lambda^m(n)_{\lambda^m} + 2i}
      &= A\{i\} \tnsr \<n\>_{\G^{(m)}}!
  \end{align*}
  From here we can get more complicated examples like
  \[
    E^{h\T}_{\psi^m(\lambda((n)_\lambda - (k)_\lambda - (n-k)_\lambda)) + 2i}
      = \Fil^i A\{i\} \tnsr {\binom nk}_{\G^{(m)}}
    % E^{t\T}_{(n)_{\lambda^m} - \lambda^m(k)_{\lambda^m} + 2i}
    %   &= A\{i\} \tnsr \left\{{\binom nk}_{\G^{(m)}}\right\}
  \]
\end{example}

\subsection{Integral representations}
\label{sub:integral}

Above, we have restricted to $\T$-representations of the form $\lambda^0$ and $\lambda^{p^i}$. For most purposes this is fine, since the divided Euler class $S^{\lambda^{p^n}} \to S^{\lambda^n}$ becomes an isomorphism after $p$-completion. However, for pinning down units we need to \emph{keep track of this isomorphism}. In particular, need this when identifying the norm map: for example, there is a norm map
\[ \pi_2 \THH \morph^N \THH^{hC_p}_{\ind^{C_p}_e 2} = \THH^{hC_p}_{(p)_\lambda} \cong \THH^{hC_p}_{\lambda^0 + (p-1)\lambda^1} \]
where $(p)_\lambda = \lambda^0 + \lambda^1 + \dotsb + \lambda^{p-1}$. The stated isomorphism multiples the $F$ map by $(p-1)!$, which must be taken into account when considering the axiom $FNx=x^p$.


We recall the following theorem of Angeltveit.

\begin{definition}[{\cite[Definition 4.5]{AngeltveitPicSpCn}}]
  Let $\ds a=(a_d)\in\prod_{1\ne d\mid n} \Z$ with each $a_d\ne0$. Let ${}^a \m W(\Z)$ denote the Mackey functor defined in the same way as the Burnside Mackey functor $\m W(\Z)$, but with restriction maps determined by
  \[ F^n_{n/k}(C_n/C_n) = \left(\prod_{d\mid k} a_d\right)C_{n/k}/C_{n/k} \]
\end{definition}

\begin{example}
  For $n=p$, the $C_p$-Mackey functor ${}^a \m W(\Z)$ is given by
  \[ \Lewis{\Z^2}{\Z}{\begin{bmatrix}a&p\end{bmatrix}}{\begin{bmatrix}0\\1\end{bmatrix}} \]
  for $a\in\{1,\dotsc,p-1\}$. This is a transversal Mackey functor, so it can equivalently be described by the data
  \[\xymatrix{
      \fade{\Z^2} \ar@[lightgray][r]^-{\fade R} \ar@[lightgray][d]_-{\fade F} \ar@{}[dr]|<<{\mbox{\fade{$\lrcorner$}}} & \Z \ar[d]^-a \\
      \Z \ar[r]_-\can & \F_p
  }\]
\end{example}

\begin{example}
  The $C_{p^2}$-Mackey functor ${}^{(a_p,a_{p^2})} \m W(\Z)$ is given by
  \[
    \Lewiss
      {\Z^3}
      {\Z^2}
      {\Z}
      {\begin{bmatrix}a_p&p&0\\0&0&p\end{bmatrix}}
      {\begin{bmatrix}a_{p^2} & p\end{bmatrix}}
      {\begin{bmatrix}0&0\\1&0\\0&1\end{bmatrix}}
      {\begin{bmatrix}0\\1\end{bmatrix}}
  \]
  for $a\in\{1,\dotsc,p-1\}$. This is a transversal Mackey functor, so it can equivalently be described by the data
  \[\xymatrix{
      \fade{\Z^2} \ar@[lightgray][r]^-{\fade R} \ar@[lightgray][d]_-{\fade F} \ar@{}[dr]|<<{\mbox{\fade{$\lrcorner$}}} & \Z \ar[d]^-a \\
      \Z \ar[r]_-\can & \F_p
  }\]
\end{example}

\begin{theorem}[{\cite[Theorem 5.1]{AngeltveitPicSpCn}}]
  Let $X$ be the $C_n$-spectrum $S^{\lambda^{n/d} - \lambda^{an/d}}$ with $d\mid n$ and $(a,d)=1$. Then $X$ is equivalent to the Mackey functor ${}^{a} \m W(\Z)$ with $a_e=\begin{cases}a&e=d\\ 1&\text{else.}\end{cases}$
\end{theorem}

\begin{definition}
  \label{dfn:prismatic-factorial}
  Let $(A, I)$ be a prism, and let $n>0$. Write $n = p^v j$ with $(p,j)=1$. We define the symbol $(n)_\prism A$ as follows: as an $A$-module, it is the ideal $I_v = I \phi(I) \dotsm \phi^{v-1}(I)$. Using the Tambara functor $\m{(A, I)}$ constructed in \cite{PrismsTambara1}, there is a natural Mackey functor structure on $\dfrac{I_v}{I_v \tnsr I_\bullet} = I_v \tnsr_A \m{(A, I)}$. The difference between $(n)_\prism A$ and $I_v$ is that the Mackey functor $\dfrac{(n)_\prism A}{(n)_\prism A \tnsr I_\bullet}$ has its $C_{p^{v+1}} \to C_{p^v}$ restriction map multiplied by $j$. To understand this formula, note that in the $q$-de Rham case we have $(n)_q = (p^v)_q (j)_{q^{p^v}}$, so $(j)_{q^{p^v}} \equiv j \in A/(q^{p^v}-1)$. Since this is natural in the prism $(A, I)$, we often just write $(n)_\prism$.

  We similarly define $(n)_{\prism^m} = (mn)_\prism \tnsr (m)_\prism^{\tnsr -1}$, $(\Phi_n)_\prism = \bigotimes_{d\mid n} {(n)_\prism ^{\tnsr\mu(n/d)}}$, $(n)_\prism! = \bigotimes_{k=1}^n {(k)_\prism}$ and $(n)_\prism? = (n-1)_\prism!$.
  % \begin{align*}
  %   (n)_{\prism^m} &= (mn)_\prism \tnsr (m)_\prism^{\tnsr -1}\\
  %   (\Phi_n)_\prism &= \bigotimes_{d\mid n} {(n)_\prism ^{\tnsr\mu(n/d)}}\\
  %   (n)_\prism! &= \bigotimes_{k=1}^n {(k)_\prism}\\
  %   (n)_\prism? &= (n-1)_\prism!
  % \end{align*}
  Over the $q$-de Rham prism, these correspond to the ideals generated by $(n)_{q^m}$, $\Phi_n(q)$, $(n)_q!$, and $\Gamma_q(n)$ respectively.
\end{definition}

\begin{proposition}
  \renewcommand{\int}{\mathrm{int}}
  Let $\alpha$ be an integral virtual $\T$-representation with irreducible decomposition $\alpha = \sum_{i=0}^\infty k_i \lambda^i$ and dimension-sequence $d_n = \dim_\C \alpha^{C_n}$ (with $d_\infty=k_0$). If $A$ is a $\delta$-ring with spherical lift $\S_A$, and $S$ is a relatively quasiregular semiperfectoid $A$-algebra, then there are canonical identifications
  \begin{align*}
    \TCmin_\alpha(S\rel \S_A)
    &= \Nyg^{\ge d_1} \prismcpl^{(1)}_{S/A}\{k_0\} \tnsr \bigotimes_{n=1}^\infty {(n)_\prism^{\tnsr k_n}}
    &
    \TP_\alpha(S\rel\S_A)
    &= \prismcpl^{(1)}_{S/A}\{k_0\} \tnsr \bigotimes_{n=1}^\infty {(n)_\prism^{\tnsr k_n}}
    \\
    &=\Nyg^{\ge d_1} \prismcpl^{(1)}_{S/A} \tnsr \bigotimes_{n=1}^\infty {(\Phi_n)_\prism^{\tnsr d_n}}
    &
    &= \prismcpl^{(1)}_{S/A} \tnsr \bigotimes_{n=1}^\infty {(\Phi_n)_\prism^{\tnsr d_n}}
  \end{align*}
\end{proposition}

\subsection{Iterated Nygaard filtrations}
\label{sub:iterated-nygaard}

\allowdisplaybreaks

The computation of $\TCmin_\rog$ and $\TP_\rog$ was simple and elegant, but somewhat boring. We can get more interesting invariants by going ``more genuine'' and considering the spectra $\TC^{(-n)} \defeq (\TR^n)^{h\T/C_{p^n}}$. To avoid confusing ourselves with Frobenius twists, we will only consider absolute prismatic cohomology; but generalizing to the relative case is just a matter of notation. (Then again, many of the advances in this paper are ``just'' notation.)

This subsection is more in the spirit of \sec\ref{sec:dRW} and can be read after that. 

\begin{definition}[Iterated Nygaard filtrations]
  \label{rmk:iterated-nygaard}
  Define \emph{higher topological negative cyclic homology} by
  \[ \TC^{-r} = (\TR^r)^{h\T/C_{p^r}} \]
  This is equivalently the pullback 
  \[ \TCmin\times_\TP \dotsb \times_\TP \TCmin, \]
  where the structure maps are $\can\colon \TCmin\to \TP$ on the right and $\varphi\colon \TCmin\to\TP$ on the left. In particular we have $\TCmin=\TC^{-0}$, while morally $\TP=\TC^{-(-1)}$. Andriopoulos studied these spectra in his thesis \cite{AndriopoulosHigherNygaard}.
\end{definition}

There are several motivations for considering these spectra:
\begin{enumerate}
  \item Nygaard's original paper \cite{Nygaard} considered not just $\Nyg^{\ge i}\dRW=\varphi^{-1}(p^i\dRW)$, but also
    \[ \Nyg^{\ge i}_r \dRW \defeq \varphi^{-r}(p^{ir}\dRW). \]
    The modern prismatic version would be
    \[ \Nyg^{\ge i}_r \prism^{(r)}_{S/A} = \varphi^{-r}(I_r^i\prism_{S/A}) \]
    where $I_r = I\phi(I)\dotsm\phi^{r-1}(I)$; in particular, the elements $\xi_r=\xi\phi^{-1}(\xi)\dotsm\phi^{-(r-1)}(\xi)$ from \cite[Lemma 3.12]{BMS1} appear naturally in this context. When $R$ is a perfectoid ring, we have
    \[ \TC^{-r}_*(R;\Z_p) = \frac{\Ainf(R)[\sigma, t_r]}{\sigma t_r - \xi_{r+1}} \quad \text{ with }|\sigma|=2,\, |t_r|=-2 \]
    by \cite[Proposition 3.4]{AndriopoulosHigherNygaard}.

  \item \label{enum:tc-r} notably, Andriopoulos introduces the idea of viewing $\TR^r$ as the quotient $\TC^{-r}/t_r$, generalizing the identification $\TCmin/t = \THH$ \cite[Proposition 3.6]{AndriopoulosHigherNygaard}. Just as $\TCmin$ is easier to understand than $\THH$ (since it has less torsion), $\TC^{-r}$ and $\TR^{h\T}$ are easier to understand than $\TR^r$ and $\TR$, and this should clarify the passage between the cyclotomic spectra perspective and the topological Cartier module/polygonic spectra perspective \cite{TCart,KMNPolygonic}.

  \item it is hard to understand what the norm
    \[ N\colon \THH^{hC_{p^n}} \to \THH^{hC_{p^{n+1}}} \]
    is really doing arithmetically, since the finite-level homotopy groups have a lot of torsion. If we could take $(-)^{h\T}$ of the norm $N\colon \THH \to \TR^1$, we might be able to get a norm $N\colon \TCmin \to \TC^{-1}$, which should be easier to understand. However, it is unclear if this works.

  \item it is natural to ask if $\RO(\T)$-grading makes sense for $\TC$. The answer turns out to be no, but you'd stumble upon $\TC^{-r}$ while trying. (However, \cite[Corollary 1.2]{DMCurves} suggests that it might be possible to define \emph{polyrepresentation}-graded $\TC$; we do not explore this in this paper.)

  \item (integral versions of) these spectra arise when constructing Habiro cohomology via refined topological Hochschild homology: see \cite[1.23]{RefinedTC-} and \cite[Theorem 6.5]{HabiroCohomology}. Relatedly, they arise in Devalapurkar-Hahn-Raksit-Yuan's work on prismatization of ring spectra \cite{PrismatizationRingSpectra} for the purpose of decompleting the Nygaard filtration, see the discussion in \cite[Chapter 7]{DevalapurkarThesis}. 

  \item Lahoti and Manam study the stacks associated to $\TR^{h\T}$ \cite[\sec1.3]{LMCohomologyModuliRingStacks}, which they call the stack of \emph{polyfiltered Cartier-Witt divisors}. A variant in positive characteristic is studied by Zhang in \cite{StackyDRW}.
\end{enumerate}

\begin{notation}
  If $\alpha=(d_0,d_1,\dotsc)$, recall that $\alpha'=\alpha^{C_p}=(d_1,d_2,\dotsc)$ and $\alpha^{(r)}=\alpha^{C_{p^r}}=(d_r,d_{r+1},\dotsc)$. To ease notation, set
  \begin{align*}
    \varphi^{-n}(I^\alpha\prismcpl_S)
      &= \Nyg^{\ge d_0,\dotsc,d_{n-1}}_n I^{d_n} \phi(I)^{d_{n+1}}\dotsm\prismcpl_S\\
      &= \left\{f \in \prismcpl_S\{d_\infty\} \,\middle|\, \varphi^k(f) \in \left(\prod_{i=0}^\infty \phi^i(I)^{d_{n-k+i}-d_\infty}\right)\prismcpl_S\{d_\infty\} \text{ for all } 0\le k\le n\right\}\\
      &= \{f\in\prismcpl_S\{d_\infty\} \mid \varphi^k(f) \in I^{\alpha^{(n-k)}}\prismcpl_S \text{ for all }0\le k\le n\}
  \end{align*}
  We also write
  \[ \~\varphi^n(\Nyg^{\ge d_0,\dotsc,d_{n-1}}) \defeq \frac{\varphi^n}{I^{d_0}\dotsm \phi^{n-1}(I)^{d_{n-1}}}(\Nyg^{\ge d_0,\dotsc,d_{n-1}}) \]
  for the image of a divided Frobenius (the map itself depends on orienting $I$, but its image does not).
\end{notation}

\begin{example}
  \label{ex:higher-tc-min}
  First we explain how to compute $\pi_\rog$ of $\TC^{-1}=(\THH^{C_p})^{h\T/C_p}$. Write $\TC^{-,\can}=\THH^{h\T}$ for the copy of $\TCmin$ that maps to $\TP$ via the canonical map, and $\TC^{-,\varphi}=(\THH^{\Phi C_p})^{h\T/C_p}$ for the copy of $\TC^-$ that maps to $\TP$ via Frobenius. Then
  \begin{align*}
    \TC^{-,\can}_\alpha(S;\Z_p) &= \varphi^{-1}(I^\alpha\prismcpl_S) = \Nyg^{\ge d_0}I^{\alpha'}\prismcpl_S\\
    \TP_\alpha(S;\Z_p) &= I^{\alpha'}\prismcpl_S\\
    \TC^{-,\varphi}_\alpha(S;\Z_p) &= \varphi^{-1}(I^{\alpha'}\prismcpl_S) = \Nyg^{\ge d_1}I^{\alpha''}\prismcpl_S
  \end{align*}
  The spectral sequence computing $\TC^{-1}_{\alpha-*}(S;\Z_p)$, where $*\in\{0,1\}$, is
  \begin{align*}
    E_0 &=\dbox{$\vcenter{\xymatrix@R=0.2em{
          \TP_\alpha(S;\Z_p) & \TC^{-,\can}_\alpha(S;\Z_p) \ar[l]\\
                       & \TC^{-,\varphi}_\alpha(S;\Z_p)
                   }}$}\\
       &=
    \dbox{$\vcenter{\xymatrix@R=0.2em{
  I^{\alpha'}\prismcpl_S & \varphi^{-1}(I^\alpha\prismcpl_S) \ar[l]_-\can\\
                          & \varphi^{-1}(I^{\alpha'}\prismcpl_S)
                                                               }}$}\\
    E_1 &=
    \dbox{$\vcenter{\xymatrix@R=0.2em{
      \dfrac{I^{\alpha'}\prismcpl_S}{\varphi^{-1}(I^\alpha)} & \\
      & \varphi^{-1}(I^{\alpha'}\prismcpl_S) \ar[ul]_-\varphi
                                                               }}$}\\
    E_\infty &=
    \dbox{$\vcenter{\xymatrix@R=0.2em{
  \dfrac{I^{\alpha'}\prismcpl_S}{\varphi^{-1}(I^\alpha) + \varphi(\varphi^{-1}(I^{\alpha'}))} & \\
      & \varphi^{-2}(I^\alpha\prismcpl_S)
                                    }}$}\\
    &=
    \dbox{$\vcenter{\xymatrix@R=0.2em{
  \dfrac{\prismcpl_S}{\Nyg^{\ge d_0} + \~\varphi(\Nyg^{\ge d_1})} \tnsr I^{\alpha'}& \\
      & \Nyg^{\ge d_0,d_1}_2 \prismcpl_S \tnsr I^{\alpha''}
    }}$}
  \end{align*}
  We conclude that
  \begin{align*}
    \TC^{-1}_\alpha(S;\Z_p)
      &= \Nyg^{\ge d_0,d_1}_2 \prismcpl_S \tnsr I^{\alpha''}\\
    \TC^{-1}_{\alpha-1}(S;\Z_p)
      &= \dfrac{\prismcpl_S}{\Nyg^{\ge d_0} + \~\varphi(\Nyg^{\ge d_1})} \tnsr I^{\alpha'}
  \end{align*}
\end{example}

For later use, it will be helpful to rewrite the second group as the iterated pushout of
\[\xymatrix{
    & \Nyg^{\ge d_0}\prismcpl_S \tnsr I^{d_1} \ar[d]_-\can \ar[r] & 0\\
    \Nyg^{\ge d_1}\prismcpl_S \ar[r]^-\varphi \ar[d] & \prismcpl_S \tnsr I^{d_1}\\
    0
}\]
(which should further be tensored with $\phi(I^{\alpha''})$, but we omit that since it is not the interesting part).

\begin{example}
  Let $R = \Z_p[\zeta_{p^\infty}]^\land_p$, and let $S=(R[x^{1/p^\infty}]/x)^\land_p$. Let $A=\Ainf(R) = \Z_p[q^{1/p^\infty}]^\land_{(p,\~\xi)}$ where $\~\xi=(p)_q$. By \cite[\sec12.2]{Prismatic}, we have
\[
\prism_S = \bigoplusp_{\nu \in \N[1/p]} A \cdot \frac{y^\nu}{\lf \nu\rf_q!}
\]
where $y=x^{1/p}$ and the completion is $(p,\~\xi)$-adic. By \cite[Lemma 12.7(1)]{Prismatic}, the Nygaard filtration is
\[
  \Nyg^{\ge i} \prism_S = \bigoplusp_\nu \xi^{\pos{i - \lf\nu\rf}} A \cdot \frac{y^\nu}{\lf\nu\rf_q!}
\]
with $\xi=\phi^{-1}(\~\xi) = (p)_{q^{1/p}}$ and $\pos x \defeq \max(x, 0)$. Using Legendre's formula to identify $\~\varphi(\Nyg^{\ge d_1}\prism_S)$, we find that
\begin{align*}
  \Nyg^{\ge d_0, d_1} \prism_S &=
  \bigoplusp_\nu \phi^{-1}(\xi)^{\pos{d_0 - \lf\nu\rf}} \xi^{\pos{d_1 - \lf\nu\rf}} A \fade{{}\cdot \frac{y^\nu}{\lf\nu\rf_q!}}\\
  \dfrac
    {\prism_S}
    {\Nyg^{\ge d_0} \prism_S + \~\varphi(\Nyg^{\ge d_1} \prism_S)}
  &=
  \bigoplusp_{d_1 \le \lf\nu\rf \le d_0} A/(\xi^{\pos{d_0 - \lf\nu\rf}}, \phi(\xi)^{\pos{\lf\nu\rf - d_1}} ) \fade{{}\cdot \frac{y^\nu}{\lf\nu\rf_q!}}
\end{align*}
\end{example}

\begin{example}
  Next we compute $\TC^{-2}_\rog(S; \Z_p)$; everything below should be tensored with $I^{\alpha^{(3)}}$, but for brevity we omit that until the end.
\begin{align*}
  &\phantom{{}\leadsto}\dbox{$\vcenter{\xymatrix@C=3em@R=2em{
  \prismcpl_S\otimes I^{d_1}\phi(I)^{d_2}
    & \Nyg^{\ge d_0}\prismcpl_S\otimes I^{d_1}\phi(I)^{d_2}
       \ar[l]_{\can}\\
  \prismcpl_S\otimes I^{d_2}
    & \Nyg^{\ge d_1}\prismcpl_S\otimes I^{d_2}
       \ar[l]_{\can}\\
  & \Nyg^{\ge d_2}\prismcpl_S
}}$}\\ 
&\leadsto
\dbox{$\vcenter{\xymatrix@C=3em@R=2em{
  \dfrac{\prismcpl_S}
        {\Nyg^{\ge d_0}\prismcpl_S} \tnsr I^{d_1}\phi(I)^{d_2} &\\
  \dfrac{\prismcpl_S}
        {\Nyg^{\ge d_1}\prismcpl_S} \tnsr I^{d_2} &\\
  & \Nyg^{\ge d_2}\prismcpl_S \ar[ul]^\varphi
}}$}\\[1em]
&\leadsto \dbox{$\vcenter{
\xymatrix@C=3em@R=2em{
  \dfrac{\prismcpl_S}
        {\Nyg^{\ge d_0}\prismcpl_S} \tnsr I^{d_1}\phi(I)^{d_2} &\\
  \dfrac{\prismcpl_S}
        {\Nyg^{\ge d_1}\prismcpl_S
          +\~\varphi(\Nyg^{\ge d_2}\prismcpl_S)} \tnsr I^{d_2} &\\
  & \Nyg^{\ge d_1,d_2}\prismcpl_S \ar@/_1em/[uul]_-{\varphi^2}
}}$}\\[1em]
&\leadsto
\dbox{$\vcenter{\xymatrix@C=3em@R=2em{
  \dfrac{\prismcpl_S}
        {\Nyg^{\ge d_0}\prismcpl_S
          +\~\varphi^2(\Nyg^{\ge d_1,d_2}\prismcpl_S)} \tnsr I^{d_1}\phi(I)^{d_2}
  \fade{{}\tnsr\phi^2(I^{\alpha^{(3)}})} &\\
  \dfrac{\prismcpl_S}
        {\Nyg^{\ge d_1}\prismcpl_S
        +\~\varphi(\Nyg^{\ge d_2}\prismcpl_S)} \tnsr I^{d_2}
  \fade{{}\otimes\phi(I^{\alpha^{(3)}})} &\\
  & \Nyg^{\ge d_0,d_1,d_2}\prismcpl_S
  \fade{{}\tnsr I^{\alpha^{(3)}}}
    }}$}
\end{align*}
The extension is the iterated pushout of
\[
\xymatrix@R=2em{
  && \Nyg^{\ge d_0}\prismcpl_S\otimes I^{d_1}\phi(I)^{d_2}
      \ar[r] \ar[d]_{\can} & 0\\
  & \Nyg^{\ge d_1}\prismcpl_S\otimes I^{d_2}
      \ar[r]^-\varphi \ar[d]_{\can} & \prismcpl_S\otimes I^{d_1}\phi(I)^{d_2}\\
  \Nyg^{\ge d_2}\prismcpl_S
      \ar[r]^-\varphi \ar[d] & \prismcpl_S\otimes I^{d_2}\\
  0 &
}
\]
\end{example}

\begin{theorem}
Let $S$ be a quasiregular semiperfectoid ring, and let $\alpha=(d_0,\ldots,d_n;d_\infty)$. Then there is a canonical identification
\[
  \TC^{-n}_{\alpha}(S;\Z_p)
  =\Nyg^{\ge d_0,\ldots,d_n}\prismcpl_S
  \fade{{}\tnsr I^{\alpha^{(n+1)}}}.
\]
The group $\TC^{-n}_{\alpha-1}(S;\Z_p)$ has a filtration with associated graded
\[
  \gr^i\TC^{-n}_{\alpha-1}(S;\Z_p)=
  \frac
    {\prismcpl_S}
    {
      \Nyg^{\ge d_{n-i}}\prismcpl_S
      +
      \~\varphi^i(\Nyg^{\ge d_{(n+1-i)}, \dotsc, d_n}\prismcpl_S)
    }
  \tnsr
  I^{\alpha^{(n+1-i)}}
\]
for $1\le i\le n$ (indexed such that $\gr^0=0$). It is given explicitly by the iterated pushout
\[
\xymatrix{
    &&& \Nyg^{\ge d_0}\prismcpl_S
    \tnsr I^{\alpha^{(1)}}
      \ar[r] \ar[d]_-\can & 0\\
                          && \Nyg^{\ge d_{(n-k-1)}}\prismcpl_S
                    \tnsr I^{\alpha^{(n-k)}}
      \ar[r]^-\varphi \ar[d]_-\can & \cdots\\
                                   & \Nyg^{\ge d_{n-k}} \prismcpl_S \tnsr I^{\alpha^{(n-k+1)}} \ar[r]^-\varphi \ar[d]_-\can & \prismcpl_S \tnsr I^{\alpha^{(n-k)}}\\
  \Nyg^{\ge d_n}\hat{\prism}_S
  \tnsr I^{\alpha^{(n+1)}}
      \ar[r]^-\varphi \ar[d] & \dots\\
  0 &
}
\]
\end{theorem}

\subsubsection{Reindexing and $\TR^{h\T}$}

To tackle $\TR^{h\T}$, it will be convenient to reindex our notation. $\TR^{h\T}$ is naturally graded on $\lim_{\alpha\mapsto\alpha'}\RO(\T)$, which we identify with bi-infinite sequences $\m{\alpha}=(\ldots,d_{-2},d_{-1},d_0,d_1,d_2,\ldots)$ such that $d_i=d_{i+1}$ for $i\gg0$. For typographical reasons, we write $\m{d}_i$ instead of $d_{-i}$. Next, we shift our indexing so that
\[
  \TCmin_{\m\alpha}(S;\Z_p)
    =\Nyg^{\ge\m d_1}\prismcpl_S\otimes I^\alpha,
  \qquad
  \TP_{\m\alpha}(S;\Z_p)
    =\prismcpl_S\otimes I^\alpha.
\]

Furthermore, we write $\prismcpl_S^{(-n)}$ for $\prismcpl_S$ considered as a $\prismcpl_S$-algebra via the map $\varphi^n$. We then write $\prismcpl_S^{(-n)}\otimes\phi^{i-n}(I)$ instead of $\prismcpl_S\otimes \phi^i(I)$, and we set
\[ I^{\m\alpha(i)} = \bigotimes_{j=-i}^\infty \phi^j(I)^{d_j}. \]
To do a computation in $\TC^{-n}_\rog$, $d_i$ in the old notation corresponds to $\m d_{(n+1-i)}$ in the new notation. Thus our previous computation of $\TC^{-1}_\rog(S;\Z_p)$ becomes
\[
\begin{aligned}
  \TC^{-1}_{\m{\alpha}}(S;\Z_p)
    &=\Nyg^{\ge\m{d}_2,\m{d}_1}\prismcpl_S^{(0)}
    \tnsr I^{\m\alpha(0)}\\
  \TC^{-1}_{\m{\alpha}-1}(S;\Z_p)
    &=\frac{\prismcpl_S^{(-1)}}
      {\Nyg^{\ge\m{d}_2}\prismcpl_S^{(-1)} +\~\varphi\left(\Nyg^{\ge\m{d}_1}\prismcpl_S^{(0)}\right)}
       \tnsr I^{\m\alpha(1)}.
\end{aligned}
\]

\begin{theorem}
Let $S$ be a quasiregular semiperfectoid ring. Then there is a canonical identification
\[
  \TC^{-n}_{\m{\alpha}}(S;\Z_p)
  =\Nyg^{\ge\m{d}_{n+1},\ldots,\m{d}_1}
  \prismcpl_S^{(0)} \tnsr I^{\m\alpha(0)}.
\]
The group $\TC^{-n}_{\m{\alpha}-1}(S;\Z_p)$ has a filtration with associated graded
\[
  \gr^i\TC^{-n}_{\m{\alpha}-1}(S;\Z_p)=
  \frac
    {\prismcpl_S^{(-i)}}
    {
      \Nyg^{\ge\m{d}_{(i+1)}}\prismcpl_S^{(-i)}
      +
      \~\varphi^i\left(\Nyg^{\ge\m{d}_i,\ldots,\m{d}_1}\prismcpl_S^{(0)}\right)
    }
  \tnsr
  I^{\m\alpha(i)}
\]
for $1\le i\le n$ (indexed such that $\gr^0=0$). It is given explicitly by the iterated pushout
\[
\xymatrix@C=4em@R=2em{
  && \Nyg^{\ge\m{d}_{n+1}}\prismcpl_S^{(-n)}
      \tnsr I^{\m\alpha(n)}
      \ar[r] \ar[d] & 0\\
  & \Nyg^{\ge\m{d}_2}\prismcpl_S^{(-1)}
      \tnsr I^{\m\alpha(1)}
      \ar[r] \ar[d] & \cdots\\
  \Nyg^{\ge\m{d}_1}\prismcpl_S^{(0)}
  \tnsr I^{\m\alpha(0)}
      \ar[r] \ar[d] & \prismcpl_S^{(-1)}
      \tnsr I^{\m\alpha(1)}\\
  0 &
}
\]
\end{theorem}

In general we do not know how to take the limit to get $\TR^{h\T}$; but see \cite[Lemma 3.23]{BMS1} for controlling $\lim\nolimits^1$ terms.

\begin{remark}
Let $\Sol_p=\lim(\cdots\xrightarrow{p}\T\xrightarrow{p}\T)$ be the $p$-adic solenoid, and let $K_i$ be the kernel of the map to the $i^\th$ copy of $\T$. If $U$ is a (possibly infinite-dimensional) representation of $\Sol_p$, say that $U$ is \emph{quasifinite-dimensional} if each $U^{K_i}$ is finite-dimensional (terminology by analogy with \cite{KMNPolygonic}). Then the grading group $\lim_{\alpha\mapsto\alpha'}\RO(\T)$ can be identified with the Grothendieck group of quasifinite-dimensional representations of $\Sol_p$.
\end{remark}
